% Source: prior_likelyhood_posterior.tex
\begin{notechapter}
% ---- chapter-local definitions ----
\forgetcs\argmax\newcommand{\argmax}{\operatorname*{arg\,max}}
% -----------------------------------
\chapter{Bayesian Reasoning}\label{ch:bayesian-reasoning}
\chaptersubtitle{Prior, likelihood and posterior}

\section{Prior, Likelihood, and Posterior}

In Bayesian reasoning, we are often interested in inferring the 
probability of a hypothesis \( \theta \) given some observed data 
\( \mathcal{D} \). This is achieved using Bayes' theorem, which 
relates the \textbf{prior}, \textbf{likelihood}, and \textbf{posterior} 
probabilities.

\subsection{Bayes' Theorem}
Bayes' theorem is stated as:
\[
p(\theta \mid \mathcal{D}) = \frac{p(\mathcal{D} \mid \theta) \, p(\theta)}{p(\mathcal{D})},
\]
where:
\begin{itemize}
    \item \( p(\theta \mid \mathcal{D}) \) is the \textbf{posterior 
          probability} of the hypothesis \( \theta \) given the data 
          \( \mathcal{D} \).
    \item \( p(\mathcal{D} \mid \theta) \) is the \textbf{likelihood} 
          of the data \( \mathcal{D} \) under the hypothesis \( \theta \).
    \item \( p(\theta) \) is the \textbf{prior probability} of the 
          hypothesis \( \theta \), representing our beliefs about 
          \( \theta \) before observing the data.
    \item \( p(\mathcal{D}) \) is the \textbf{marginal likelihood} 
          (or evidence), which acts as a normalizing constant to ensure 
          the posterior probability sums to 1.
\end{itemize}

\subsection{Prior Probability}
The prior probability \( p(\theta) \) encodes our initial beliefs 
about the hypothesis \( \theta \) before observing any data. For 
example, if \( \theta \) represents the probability of a coin landing 
heads, a reasonable prior might be a uniform distribution, reflecting 
no initial bias toward heads or tails.

\subsection{Likelihood}
The likelihood \( p(\mathcal{D} \mid \theta) \) quantifies how well 
the hypothesis \( \theta \) explains the observed data \( \mathcal{D} \). 
For example, if \( \theta \) is the probability of a coin landing 
heads, and we observe 7 heads in 10 flips, the likelihood measures 
how probable this outcome is under the hypothesis \( \theta \).

\subsection{Posterior Probability}
The posterior probability \( p(\theta \mid \mathcal{D}) \) represents 
our updated beliefs about the hypothesis \( \theta \) after observing 
the data \( \mathcal{D} \). It combines the prior probability and the 
likelihood using Bayes' theorem. The posterior is proportional to the 
product of the prior and the likelihood:
\[
p(\theta \mid \mathcal{D}) \propto p(\mathcal{D} \mid \theta) \, p(\theta).
\]

\subsection{Marginal Likelihood}
The marginal likelihood \( p(\mathcal{D}) \) is computed by summing 
(or integrating) over all possible hypotheses:
\[
p(\mathcal{D}) = \sum_{\theta} p(\mathcal{D} \mid \theta) \, p(\theta)
\quad\text{or}\quad
p(\mathcal{D}) = \int p(\mathcal{D} \mid \theta) \, p(\theta)\, d\theta .
\]
This ensures that the posterior probability is properly normalized.

\subsection{Generative Model and Posterior Inference}
Bayesian reasoning allows us to infer the posterior distribution 
\( p(\theta \mid \mathcal{D}) \) of a variable \( \theta \) given 
observed data \( \mathcal{D} \). This is achieved by combining a 
\textbf{forward} or \textbf{generative model} \( p(\mathcal{D} \mid \theta) \) 
with a \textbf{prior belief} \( p(\theta) \) about the variable. The 
generative model describes how the data \( \mathcal{D} \) is generated 
from the variable \( \theta \), while the prior encodes our initial 
beliefs about plausible values of \( \theta \).

\subsection{Maximum a Posteriori (MAP) Estimation}
The \textbf{most probable a posteriori (MAP)} setting is the value of 
\( \theta \) that maximizes the posterior distribution:
\[
\theta^* = \argmax_{\theta} p(\theta \mid \mathcal{D}).
\]
This represents the most plausible value of \( \theta \) given the 
observed data \( \mathcal{D} \) and the prior \( p(\theta) \).

\subsection{Connection to Maximum Likelihood (ML)}
When the prior \( p(\theta) \) is \textbf{flat} (i.e., a constant that 
does not depend on \( \theta \)), the posterior distribution is 
proportional to the likelihood:
\[
p(\theta \mid \mathcal{D}) \propto p(\mathcal{D} \mid \theta).
\]
In this case, the MAP solution reduces to the \textbf{Maximum Likelihood 
(ML)} estimate:
\[
\theta^* = \argmax_{\theta} p(\mathcal{D} \mid \theta).
\]
The ML estimate is the value of \( \theta \) that maximizes the 
likelihood of the observed data \( \mathcal{D} \), ignoring any prior 
information about \( \theta \).

\subsection{Example: Coin Tossing}
Consider a simple example of inferring the bias \( \theta \) of a coin 
from observed data. Suppose we flip a coin \( N \) times and observe 
\( h \) heads. The likelihood of the observed sequence of flips is
\[
p(\mathcal{D} \mid \theta) = \theta^h (1 - \theta)^{N - h}.
\]
(If only the number of heads is recorded, the binomial likelihood has an extra factor \( \binom{N}{h} \), which does not depend on \( \theta \) and therefore does not change the posterior.)
If we assume a uniform prior \( p(\theta) = 1 \) for \( \theta \in [0, 1] \), 
the posterior distribution is:
\[
p(\theta \mid \mathcal{D}) \propto \theta^h (1 - \theta)^{N - h}.
\]
This is a Beta distribution, \( \text{Beta}(\theta \mid h + 1, N - h + 1) \), 
which reflects our updated beliefs about the coin's bias after observing 
the data.

\subsection{Example: Pendulum}
Consider an idealised pendulum where \( x_t \) is the angular displacement 
of the pendulum at time \( t \). Assuming that the measurements are 
independent given the parameter \( \theta \), the likelihood of a sequence 
of observations \( x_1, \dots, x_T \) is given by:
\[
p(x_1, \dots, x_T \mid \theta) = \prod_{t=1}^T p(x_t \mid \theta).
\]
If the model is perfect and the measurements are perfect, the physical 
model is:
\[
x_t = \sin(\theta t),
\]
where \( \theta \) represents the unknown physical constants of the 
pendulum (e.g., \( \theta = \sqrt{g/L} \), where \( g \) is the 
gravitational acceleration and \( L \) is the length of the pendulum). 

However, if the measurement instrument is poor, with a known variance 
\( \sigma^2 \), the model becomes:
\[
x_t = \sin(\theta t) + \epsilon_t,
\]
where \( \epsilon_t \) is zero-mean Gaussian noise with variance \( \sigma^2 \).

We can place a prior \( p(\theta) \) over the possible values of \( \theta \), 
expressing our prior belief in the appropriateness of different values. 
The posterior distribution is then given by:
\[
p(\theta \mid x_1, \dots, x_T) \propto p(\theta) \prod_{t=1}^T p(x_t \mid \theta).
\]
Despite noisy measurements, the posterior over the possible values of \( \theta \) 
becomes strongly peaked, as shown in Figure~1.1 of D.~Barber, \emph{Bayesian Reasoning and Machine Learning}.

\subsubsection{Derivation of the Posterior}
The expression for the posterior distribution:
\[
p(\theta \mid x_1, \dots, x_T) \propto p(\theta) \prod_{t=1}^T p(x_t \mid \theta)
\]
is derived from \textbf{Bayes' theorem} and the assumption that the observations 
\( x_1, \dots, x_T \) are \textbf{conditionally independent} given \( \theta \).

\paragraph{Bayes' Theorem}
Bayes' theorem states that the posterior distribution is proportional to the 
product of the prior and the likelihood:
\[
p(\theta \mid x_1, \dots, x_T) = \frac{p(x_1, \dots, x_T \mid \theta) \, p(\theta)}{p(x_1, \dots, x_T)}.
\]
Since the marginal likelihood \( p(x_1, \dots, x_T) \) does not depend on \( \theta \), 
we can write:
\[
p(\theta \mid x_1, \dots, x_T) \propto p(x_1, \dots, x_T \mid \theta) \, p(\theta).
\]

\paragraph{Conditional Independence}
The key assumption is that the observations \( x_1, \dots, x_T \) are conditionally 
independent given \( \theta \). This means:
\[
p(x_1, \dots, x_T \mid \theta) = \prod_{t=1}^T p(x_t \mid \theta).
\]
Substituting this into Bayes' theorem gives:
\[
p(\theta \mid x_1, \dots, x_T) \propto p(\theta) \prod_{t=1}^T p(x_t \mid \theta).
\]

\paragraph{Interpretation}
\begin{itemize}
    \item The \textbf{prior} \( p(\theta) \) encodes our initial beliefs about \( \theta \).
    \item The \textbf{likelihood} \( \prod_{t=1}^T p(x_t \mid \theta) \) measures how well 
          \( \theta \) explains the observed data.
    \item The \textbf{posterior} combines the prior and likelihood to provide updated 
          beliefs about \( \theta \).
\end{itemize}

\subsection{Summary}
Bayesian reasoning provides a principled framework for updating beliefs 
in light of new evidence. The prior represents initial beliefs, the 
likelihood measures how well the hypothesis explains the data, and the 
posterior combines these to provide updated beliefs. The marginal 
likelihood ensures proper normalization of the posterior. Additionally, 
the MAP estimate provides the most plausible value of \( \theta \) given 
the data and prior, while the ML estimate maximizes the likelihood under 
a flat prior.
\end{notechapter}
